% ==============================================================================
% SLIDES - SEMANA 10: GESTÃO DE RISCOS, COMPLIANCE & SEGURANÇA DA INFORMAÇÃO
% ==============================================================================

% 1. CAPA
\senaicover
  {Riscos, Compliance \& Segurança}
  {Gestão de Riscos (ISO 31000), Matriz de Riscos 5x5, LGPD e Entrega Parcial 2}
  {Unidade Curricular: Governança de TI}
  {Prof. André Souza}

% 2. AGENDA
\begin{senaiframe}{Visão Geral \& Agenda}{Estrutura da Aula}
  \begin{columns}[t]
    \begin{column}{0.48\textwidth}
      \begin{tcolorbox}[senaibluecard, title={1. Gestão de Riscos}]
        \begin{itemize}
          \item Conceito de Risco de TI (Ameaça x Vulnerabilidade).
          \item As 4 etapas da ISO 31000.
        \end{itemize}
      \end{tcolorbox}
      \vspace{4pt}
      \begin{tcolorbox}[senaiambercard, title={3. Compliance Regulatório}]
        \begin{itemize}
          \item Requisitos da LGPD (Lei 13.709/2018).
          \item ISO/IEC 27001 e a Tríade CID.
        \end{itemize}
      \end{tcolorbox}
    \end{column}
    \begin{column}{0.48\textwidth}
      \begin{tcolorbox}[senairedcard, title={2. Matriz de Riscos}]
        \begin{itemize}
          \item Probabilidade x Impacto (Matriz 5x5).
          \item As 4 respostas: Mitigar, Evitar, Transferir, Aceitar.
        \end{itemize}
      \end{tcolorbox}
      \vspace{4pt}
      \begin{tcolorbox}[senaigreencard, title={4. Entrega Parcial 2}]
        \begin{itemize}
          \item Submissão dos Capítulos 5 e 6.
          \item Segundo marco do projeto PDGTI.
        \end{itemize}
      \end{tcolorbox}
    \end{column}
  \end{columns}
\end{senaiframe}

% 3. CONCEITO DE RISCO DE TI
\begin{senaiframe}{Teoria • Módulo 1}{O que é Risco de TI? (ISO 31000)}
  \begin{columns}[t]
    \begin{column}{0.48\textwidth}
      \begin{tcolorbox}[senaibluecard, title={Definição Geral}]
        \begin{itemize}
          \item Efeito da incerteza nos objetivos da organização.
          \item Probabilidade de um evento afetar negativamente os negócios.
        \end{itemize}
      \end{tcolorbox}
    \end{column}
    \begin{column}{0.48\textwidth}
      \begin{tcolorbox}[senaicard, title={Contexto da TI}]
        \begin{itemize}
          \item Falhas de sistemas, ataques virtuais ou erros operacionais.
          \item Danos financeiros, operacionais e reputacionais.
        \end{itemize}
      \end{tcolorbox}
    \end{column}
  \end{columns}
\end{senaiframe}

% 4. AMEAÇA VS VULNERABILIDADE
\begin{senaiframe}{Teoria • Módulo 1}{Ameaça vs. Vulnerabilidade vs. Impacto}
  \begin{columns}[t]
    \begin{column}{0.48\textwidth}
      \begin{tcolorbox}[senairedcard, title={Ameaça \& Vulnerabilidade}]
        \begin{itemize}
          \item \textbf{Ameaça:} Agente externo/interno (ex: hacker, enchente).
          \item \textbf{Vulnerabilidade:} Fraqueza no sistema (ex: senha fraca).
        \end{itemize}
      \end{tcolorbox}
    \end{column}
    \begin{column}{0.48\textwidth}
      \begin{tcolorbox}[senaiambercard, title={Equação do Risco}]
        \begin{itemize}
          \item \textbf{Risco = Ameaça x Vulnerabilidade x Impacto}.
          \item Sem vulnerabilidade, a ameaça não se concretiza.
        \end{itemize}
      \end{tcolorbox}
    \end{column}
  \end{columns}
\end{senaiframe}

% 5. ETAPAS DA GESTÃO DE RISCOS
\begin{senaiframe}{Teoria • Módulo 1}{As 4 Etapas da Gestão de Riscos}
  \begin{columns}[t]
    \begin{column}{0.48\textwidth}
      \begin{tcolorbox}[senaibluecard, title={Identificação \& Análise}]
        \begin{itemize}
          \item \textbf{1. Identificação:} Mapeamento das ameaças aos ativos.
          \item \textbf{2. Análise:} Determinar probabilidade e impacto.
        \end{itemize}
      \end{tcolorbox}
    \end{column}
    \begin{column}{0.48\textwidth}
      \begin{tcolorbox}[senaigreencard, title={Avaliação \& Resposta}]
        \begin{itemize}
          \item \textbf{3. Avaliação:} Priorização na matriz de risco.
          \item \textbf{4. Resposta:} Definir plano de ação/mitigação.
        \end{itemize}
      \end{tcolorbox}
    \end{column}
  \end{columns}
\end{senaiframe}

% 6. MATRIZ DE RISCOS 5X5
\begin{senaiframe}{Teoria • Módulo 2}{A Matriz de Riscos 5x5}
  \begin{columns}[t]
    \begin{column}{0.48\textwidth}
      \begin{tcolorbox}[senairedcard, title={Eixos de Avaliação}]
        \begin{itemize}
          \item \textbf{Probabilidade (1 a 5):} Rara até Quase Certa.
          \item \textbf{Impacto (1 a 5):} Insignificante até Catastrófico.
        \end{itemize}
      \end{tcolorbox}
    \end{column}
    \begin{column}{0.48\textwidth}
      \begin{tcolorbox}[senaiambercard, title={Categorias de Severidade}]
        \begin{itemize}
          \item \textbf{Baixo, Médio, Alto e Crítico}.
          \item Foco prioritário de investimento nos riscos Críticos.
        \end{itemize}
      \end{tcolorbox}
    \end{column}
  \end{columns}
\end{senaiframe}

% 7. RESPOSTAS: MITIGAR E EVITAR
\begin{senaiframe}{Teoria • Módulo 2}{Estratégias: Mitigar \& Evitar}
  \begin{columns}[t]
    \begin{column}{0.48\textwidth}
      \begin{tcolorbox}[senaibluecard, title={Mitigar (Reduzir)}]
        \begin{itemize}
          \item Implementar controles técnicos para diminuir risco.
          \item *Exemplo:* Instalação de MFA e criptografia.
        \end{itemize}
      \end{tcolorbox}
    \end{column}
    \begin{column}{0.48\textwidth}
      \begin{tcolorbox}[senairedcard, title={Evitar (Eliminar)}]
        \begin{itemize}
          \item Descontinuar tecnologia ou atividade arriscada.
          \item *Exemplo:* Desligar servidores sem patches de segurança.
        \end{itemize}
      \end{tcolorbox}
    \end{column}
  \end{columns}
\end{senaiframe}

% 8. RESPOSTAS: TRANSFERIR E ACEITAR
\begin{senaiframe}{Teoria • Módulo 2}{Estratégias: Transferir \& Aceitar}
  \begin{columns}[t]
    \begin{column}{0.48\textwidth}
      \begin{tcolorbox}[senaiambercard, title={Transferir (Compartilhar)}]
        \begin{itemize}
          \item Repassar o impacto financeiro a terceiros.
          \item *Exemplo:* Seguro cibernético ou contratação de nuvem.
        \end{itemize}
      \end{tcolorbox}
    \end{column}
    \begin{column}{0.48\textwidth}
      \begin{tcolorbox}[senaicard, title={Aceitar (Risco Residual)}]
        \begin{itemize}
          \item Conviver com o risco quando o custo do controle supera a perda.
          \item Decisão formalizada e aprovada pela diretoria.
        \end{itemize}
      \end{tcolorbox}
    \end{column}
  \end{columns}
\end{senaiframe}

% 9. UNIVERSIO DE COMPLIANCE
\begin{senaiframe}{Teoria • Módulo 3}{O Universo de Compliance Regulatório}
  \begin{columns}[t]
    \begin{column}{0.48\textwidth}
      \begin{tcolorbox}[senaibluecard, title={Regulação da TI}]
        \begin{itemize}
          \item Exigências legais sobre armazenamento de dados e privacidade.
          \item Auditabilidade de processos operacionais de TI.
        \end{itemize}
      \end{tcolorbox}
    \end{column}
    \begin{column}{0.48\textwidth}
      \begin{tcolorbox}[senairedcard, title={Sanções de Não-Conformidade}]
        \begin{itemize}
          \item Multas milionárias, paralisação de sistemas e dano à marca.
        \end{itemize}
      \end{tcolorbox}
    \end{column}
  \end{columns}
\end{senaiframe}

% 10. A LGPD NA TI
\begin{senaiframe}{Teoria • Módulo 3}{A LGPD (Lei 13.709/2018) na TI}
  \begin{columns}[t]
    \begin{column}{0.48\textwidth}
      \begin{tcolorbox}[senaiambercard, title={Diretrizes Gerais}]
        \begin{itemize}
          \item Proteção de dados pessoais e privacidade dos titulares.
          \item Multas de até 2\% do faturamento (limite de R$ 50 mi).
        \end{itemize}
      \end{tcolorbox}
    \end{column}
    \begin{column}{0.48\textwidth}
      \begin{tcolorbox}[senaigreencard, title={Controles de TI}]
        \begin{itemize}
          \item Anonimização de dados, gestão de consentimento e DPO.
          \item Relatório de Impacto à Proteção de Dados (RIPD).
        \end{itemize}
      \end{tcolorbox}
    \end{column}
  \end{columns}
\end{senaiframe}

% 11. ISO 27001 E TRÍADE CID
\begin{senaiframe}{Teoria • Módulo 3}{ISO/IEC 27001 \& Tríade CID}
  \begin{columns}[t]
    \begin{column}{0.48\textwidth}
      \begin{tcolorbox}[senaibluecard, title={A Tríade CID}]
        \begin{itemize}
          \item \textbf{Confidencialidade:} Acesso restrito a pessoas autorizadas.
          \item \textbf{Integridade:} Dados exatos e sem modificações indevidas.
          \item \textbf{Disponibilidade:} Acesso garantido quando necessário.
        \end{itemize}
      \end{tcolorbox}
    \end{column}
    \begin{column}{0.48\textwidth}
      \begin{tcolorbox}[senaicard, title={SGSI}]
        \begin{itemize}
          \item Sistema de Gestão de Segurança da Informação formal.
        \end{itemize}
      \end{tcolorbox}
    \end{column}
  \end{columns}
\end{senaiframe}

% 12. OUTROS MARCOS
\begin{senaiframe}{Teoria • Módulo 3}{Outros Marcos Regulatórios (PCI-DSS, BACEN)}
  \begin{columns}[t]
    \begin{column}{0.48\textwidth}
      \begin{tcolorbox}[senaibluecard, title={PCI-DSS}]
        \begin{itemize}
          \item Padrão para transações com cartões de crédito.
        \end{itemize}
      \end{tcolorbox}
    \end{column}
    \begin{column}{0.48\textwidth}
      \begin{tcolorbox}[senaicard, title={BACEN (Bancos)}]
        \begin{itemize}
          \item Resoluções sobre cibersegurança no setor financeiro.
        \end{itemize}
      \end{tcolorbox}
    \end{column}
  \end{columns}
\end{senaiframe}

% 13. ESTUDO DE CASO PRÁTICO
\begin{senaiframe}{Estudo de Caso}{Rede Hospitalar: Sequestro por Ransomware}
  \begin{columns}[t]
    \begin{column}{0.48\textwidth}
      \begin{tcolorbox}[senairedcard, title={Incidente Real}]
        \begin{itemize}
          \item Criptografia do banco de exames; resgate de US$ 2 milhões.
          \item Causa: Servidores desatualizados e sem backup offline.
        \end{itemize}
      \end{tcolorbox}
    \end{column}
    \begin{column}{0.48\textwidth}
      \begin{tcolorbox}[senaigreencard, title={Resposta de Governança}]
        \begin{itemize}
          \item Adotado backup imutável com testes DRP mensais.
          \item Microssegmentação da rede hospitalar.
        \end{itemize}
      \end{tcolorbox}
    \end{column}
  \end{columns}
\end{senaiframe}

% 14. REFERÊNCIAS BIBLIOGRÁFICAS
\begin{senaiframe}{Referências Bibliográficas}{Fontes \& Leitura Recomendada}
  \begin{tcolorbox}[senaicard, title={Obras de Referência}]
    \begin{itemize}
      \item \textbf{FREITAS, Daniel Paulo Paiva.} *Proteção e governança de dados.* Curitiba: Contentus, 2020.
      \item \textbf{ANDRADE, Alexandre Francisco.} *Gestão de Compliance.* Curitiba: Contentus, 2020.
      \item \textbf{ISO/IEC.} *ISO/IEC 27001: Information security, cybersecurity and privacy protection.* ISO, 2022.
    \end{itemize}
  \end{tcolorbox}
\end{senaiframe}

% 15. APLICAÇÃO NO PDGTI
\begin{senaiframe}{Aplicação no PDGTI}{Orientações para a ENTREGA PARCIAL 2}
  \vspace{0.2cm}
  \begin{tcolorbox}[senaigreencard, title={Entrega Parcial 2 (Capítulos 5 e 6) – Data: 09/Out}]
    \begin{itemize}
      \item \textbf{Capítulo 5:} Indicadores de Desempenho (KPIs Negócio/TI).
      \item \textbf{Capítulo 6:} Matriz de Riscos 5x5, Mitigação e LGPD.
      \item Consolidação das correções solicitadas na Entrega Parcial 1.
    \end{itemize}
  \end{tcolorbox}
\end{senaiframe}
